\documentclass[letter]{article}

\usepackage{listings,geometry,amsmath}

\lstset{language=Matlab,basicstyle=\small\ttfamily}
\geometry{margin=1in,nohead}

\setlength{\parindent}{0em}
\setlength{\parskip}{1.5ex plus0.5ex minus0.5ex}
\renewcommand{\baselinestretch}{1.2}

\begin{document}
\large{ERTH 467: Lab 00}\hfill \large{September 4, 2026}\\

\section*{Introduction}
This first lab is simply to introduce you to the software that we will
use for the remainder of the labs. You don't need to hand in a formal lab write-up for this lab.  However, you should complete the exercises and make sure that you understand the
material.  If you have any questions, please ask.  

We use the Matlab software package.  Matlab is a really easy to use programming language with
tons of built-in functions, including a bundle of plotting
routines. If you already have some programming experience, this lab
will be easy.  If you have never programmed before, then this lab will
provide a gentle introduction.

If you know Python, notice that Matlab syntax is very similar to Python's NumPy library.  
If you know C, notice that Matlab syntax is very similar to C syntax.
If you know neither, then this lab will provide a gentle introduction.

The reason for using Matlab instead of numpy is that Matlab is clearer for signal processing.
A large number of signal processing functions are built into Matlab, and the syntax is more clear than numpy.  
For example, the Fourier transform in Matlab is `fft' while in numpy it is `numpy.fft.fft'.  
The Matlab syntax is much easier to read.

Another reason for using Matlab, is that occassionally we will the CREWES toolbox, 
which is a collection of Matlab functions for seismic processing. 
Many of the functions we will develop in this course have some similar version in the CREWES toolbox.
You can find the CREWES toolbox at: `http://www.crewes.org/ResearchLinks/Software/MatlabToolbox/'.

Similarly, Matlab has a large number of built-in functions for signal processing,
including the Fourier transform, filtering, and plotting.  You can find the Matlab documentation at:
`http://www.mathworks.com/help/matlab/'.

NOTE: you can use the crewes toolbox to check your results or get ideas for your own code, 
but you should write your own code for the lab exercises.

On the other hand, if you prefer to use Python, please feel free to do so.  
The lab exercises can be completed in either language.  
If you choose to use Python, please make sure that you have the numpy and matplotlib libraries installed.


\section*{Manipulating matrices, vectors and scalars}
Matlab stores information in \emph{variables}.  Almost every variable in
Matlab is stored as a \emph{matrix}.  Try entering the following
commands (without the bracketed and italicized comments, or the
command prompt, $>>$)\ldots
\begin{enumerate}
\small
\item $>>$ a = 2 \emph{(assigns 2 to $a$)}
\item $>>$ x = [1;2;3] \emph{(assigns a column vector to $x$)}
\item $>>$ y = [4,5,6] \emph{(assigns a row vector to $y$)}
\item $>>$ z = [0:2:20] \emph{(assigns a row vector
$\left[\begin{array}{ccccc} 0 & 2 & 4 & \cdots & 20 \end{array}\right]$ to $z$)}
\item $>>$ A = [1,2,3;4,5,6;7,8,0] \emph{(assigns a matrix to $A$)}
\item $>>$ x \emph{(prints $x$)}
\item $>>$ whos \emph{(prints a table of stored variables)}
\item $>>$ size(A) \emph{(prints the size of A)}
\item $>>$ size(z) \emph{(prints the size of z)}
\item $>>$ length(z) \emph{(prints the length of z)}
\end{enumerate}
`a' is a scalar ($1\times 1$ matrix); `x' is a vector ($3\times 1$
matrix); `y' is a vector ($1\times 3$ matrix); `A' is a matrix ($3\times 3$
matrix). $>>$ is called the command prompt.

You can access subsets of the matrices and vectors.  Try the
following\ldots
\begin{enumerate}
\small
\item $>>$ x(2) \emph{(access the 2nd element of the vector)}
\item $>>$ x(2:3) \emph{(access the 2nd through 3rd elements of the vector)}
\item $>>$ A(2,3) \emph{(access the element, $A_{23}$ of the matrix)}
\item $>>$ A(1:2,2:3) \emph{(access the elements $A_{ij}$ for
$i=1\ldots 2, j=2\ldots3$)} 
\end{enumerate}

You can stitch vectors and matrices together in any logical
order.  For example, try the following\ldots
\begin{enumerate}
\small
\item $>>$ [x',y] \emph{(attach $y$ to the tail of $x^T$)}
\item $>>$ [A,x] \emph{(attach a column to the right side of $\mathbf{A}$)}
\end{enumerate}
\section*{Some Linear Algebra}
Matlab supports pretty much every operation that you can find in your
linear algebra text-book (Here we look at just a few simple
examples).  Try the following\ldots
\begin{enumerate}
\small
\item $>>$ A' \emph{(returns the transpose of $A$)}
\item $>>$ A*x \emph{(matrix-vector multiplication of $A$ and $x$)}
\item $>>$ x+y' \emph{(vector addition of $x$ and $y^T$)}
\item $>>$ x.*y' \emph{(element-wise multiplication of $x$ and $y^T$)}
\item $>>$ x./y' \emph{(element-wise division of $x$ and $y^T$)}
\item $>>$ a*x \emph{(multiplication of a vector by a scalar)}
\item $>>$ x+a \emph{(vector addition of $x$ and a vector of the same
size with elements $a$)}
\end{enumerate}
Of course this is just a short list.  If you have time, you can look
at the on-line documentation for more features.

\section*{Plotting}
Matlab has some easy-to-use and useful plotting routines.  In this lab we will
use the `plot' routine.  To learn about other plotting functions in
Matlab you can use the `help' command.  Try the following...
\begin{enumerate}
\small
\item $>>$ help plot
\item $>>$ help stem
\item $>>$ help subplot
\end{enumerate}
You don't have to read the instructions\ldots Just keep in mind that
you will have to use these plotting routines for every lab exercise,
including this one.

\section*{Some Computer Science}
Matlab is a fully featured \emph{scripting} language supporting functions,
loops, and conditional statements.  This means that commands can be
typed into a text file which can be interpreted by Matlab.

\begin{itemize}
\item open a new text file in a text editor called
  `my\_first\_script.m'
\item In Matlab change directories to the one that contains your
  file by typing, at the Matlab prompt, `cd path/to/file'.  This will
  ensure that Matlab can see your text file.  If you are not familiar
  with the `cd' command, please ask.
\end{itemize}

We will enter commands into this text file.  For this first example,
our objective is very simple:

\begin{itemize}
\item Write some computer code that will produce two vectors,
  $\mathbf{x}$ and $\mathbf{y}$ such that
  $\mathbf{y}=\sin(\mathbf{x})$. Instructions 
  follow\ldots
\end{itemize}

%\vfill
%\hfill \emph{continued on next page...}

%\newpage

\subsection*{looping}

enter the following text (code) into `my\_first\_script.m' \\

\lstinputlisting[frame=tb]{my_first_script.m}

`ii' is the loop index; it is set from the elements of the row vector:
$\left[\begin{array}{ccccc} 1 & 2 & \cdots & 100 \end{array}\right]$.
Matlab \emph{loops} through the elements of this vector, executing the
contents of the loop (the lines between `for ii=1:1:100' and `end')
for each value of `ii'.  You can run this code by entering the file
name (without the suffix) at the command
prompt.  Make sure you try this to see if your code works.  You should
see a window pop up with a plot of your curve (exciting, eh?).

Alternatively, you can use the `while' looping command.  To
learn more about the `while' looping command, use the already
mentioned Matlab help feature.

\begin{itemize}
\item Try editing your `my\_first\_script.m' file, replacing the `for' loop
with a `while' loop.
\end{itemize}

Finally, it should be noted that the above code can be written without
using a loop.  An idea which you will need shortly.

\subsection*{conditionals}
Create a second file called `my\_second\_script.m' and add the
following code.  The `if' statement is the conditional.  If the
condition `x(ii)$<$0' is true, the statement, `y(ii)=-x(ii)' is
executed.  Otherwise, `y(ii)=x(ii)' is executed.

\lstinputlisting[frame=tb]{my_second_script.m}

Try running this new script (wow, even more exiting!).

\subsection*{writing functions}
So far we have written two Matlab scripts.  But, what if we want to
use the script over and over again with different values for $n$, $T$
and so on.  This is what functions are for.  Each function must be in
its own file.  The first line of code in the function file will always
look like,\\

\centerline{function[return\_val1,return\_val2,\ldots] =
my\_function(arg1,arg2,\ldots)}\vspace{.5cm}

This line is called the function header. The function's name is
`my\_function', and the function must be written in a text
file called `my\_function.m'.  The remainder of the code in the
function file will use the variables: `arg1', `arg2', \ldots; will
assign values to variables: `return\_val1', `return\_val2, \ldots; and
will look similar to the code in the scripts that you have already written.

\section*{Exercise}
In this exercise you will write a function for the code in
`my\_second\_script.m'.  Once again, your goal is to write code that
computes the vectors $\mathbf{x}$ and $\mathbf{y}$ such that,
\begin{equation}
  y=f(x)=
  \left\{
  \begin{array}{r@{\quad,\quad}l}
    -x & x < 0\\
    x  & x \geq 0
  \end{array}
  \right.
  \label{eqn:saw}
\end{equation}

\begin{itemize}
\item Open a new file called `saw\_tooth.m'
\item Use the following function header:
\end{itemize}

\centerline{function[x,y] = saw\_tooth(n,x0)}
\vspace{.5cm}

\begin{itemize}
\item The function should NOT plot $y(x)$.  It should only compute $y$
  and $x$.
\end{itemize}

The body of the function is for you to write.  Of course it will look
very similar to `my\_second\_script.m'; however, I am placing the
following restriction on your code:\\

\centerline{YOU MAY NOT USE LOOPS OR AN INDEX VARIABLE SUCH AS `ii'!}
\vspace{0.5cm}

Once you've written your function, make sure it works!  Try using
it.  For example, try:

\begin{enumerate}
\small
\item $>>$ [x,y] = saw\_tooth(101,-2);
\item $>>$ plot(x,y)
\end{enumerate}

Finally, let's put your function to use inside a Matlab script.  In
this script you will plot your function and an approximation to that
function using sine and cosine functions (ie. a Fourier
series) such that $f(x)$ in equation (\ref{eqn:saw}) is approximated
as,
\begin{displaymath}
  f(x) \approx \frac{a_\circ}{2} + \sum_{m=1}^{M}\left[a_m\sin\left(\frac{m\pi
  x}{2}\right) + b_m\cos\left(\frac{m\pi x}{2}\right)\right]
\end{displaymath}
where
\begin{displaymath}
\begin{array}{ll}
a_\circ = 2 \\
a_m = 8/\left(m\pi\right)^2 & \text{for odd } m \\
a_m = 0 & \text{for even } m \\
b_m = 0 & \text{for } m=1,2,\ldots
\end{array}
\end{displaymath}
and $f(x)$ is approximated and plotted for $M=1$ through $M=9$.
Detailed instructions follow\ldots

\begin{itemize}
\item open a new file called `fourier.m'.
\end{itemize}

Input the following code into the text file.

\lstinputlisting[frame=tb]{fourier.m}

Run the script.  If your typing skills are sufficient, you should
see something interesting.  Namely that the Fourier series works!

%\section*{HAND IN}
%DO NOT HAND IN A FORMAL LAB WRITE UP!!!  Do hand in\ldots
%\begin{enumerate}
%\item A print-out of your function, `saw\_tooth.m' (with your name on it).
%\item A print-out of the plot produced by the script, `fourier.m'.
%\item Add comments to `fourier.m' explaining what each line does.  One
%  sentence per line of code is more than enough.  Use your best judgment.
%  Hand in a print-out of your commented  code.
%\end{enumerate}

You can find some useful Matlab tutorials at:
`http://www.math.siu.edu/matlab/tutorials.html'


\end{document}

